\section{Monge Problem between Measures}
\label{sec-monge}
\subsection{Measures}
\label{sec-measures}
\paragraph{Histograms.}
\begin{defn}[Probability simplex]\label{def-probability-simplex}
The probability simplex of length $n$ is \(\simplex_n \eqdef \enscond{\a \in \RR_+^n}{ \sum_{i=1}^n \a_i = 1 }.\)
Its elements are also called probability vectors or histograms.
\end{defn}
\paragraph{Discrete measure, empirical measure.}
\begin{defn}[Discrete measure]\label{def-discrete-measure}
A discrete measure with weights $\a$ and locations $x_1,\dots,x_n\in\X$ is
\eql{\label{eq-discr-meas}
\al = \sum_{i=1}^n \a_i \de_{x_i},
}
where $\de_x$ is the Dirac mass at position $x$. It is a probability measure if $\a\in\simplex_n$, and a positive measure if all weights $\a_i$ are nonnegative.
\end{defn}
\paragraph{General measures.}
A Dirac measure $\de_x$ is then defined as $\de_x(A)=1$ if $x \in A$ and $0$ otherwise, and this extends by linearity for discrete measures of the form~\eqref{eq-discr-meas} as

\(\al(A) = \sum_{x_i \in A} \a_i\)
\paragraph{Polish metric spaces.}
\begin{defn}[Polish metric space]\label{def-polish-metric-space}
A metric space $(\X,d)$ is Polish if it is complete and separable: every Cauchy sequence converges in $\X$, and $\X$ contains a countable dense subset. More generally, a topological space is called Polish if its topology can be induced by some complete separable metric.
\end{defn}
\(\Ss=C([0,1];\RR^d),\qquad d_\infty(\gamma,\eta) \eqdef \sup_{t\in[0,1]}\norm{\gamma(t)-\eta(t)}\)
\begin{defn}[Support of a measure]\label{def:support}
Let $\al$ be a positive Borel measure on a metric space $\Xx$. Its support is \(\supp(\al) \eqdef \enscond{x\in\Xx}{\al(B(x,r))>0\text{ for every }r>0}.\)
If $\Xx$ is separable, this is the smallest closed set of full $\al$-mass.
\end{defn}
\paragraph{Radon measures.}
\begin{defn}[Radon and probability measures]\label{def-radon-probability-measures}
A locally finite positive Borel measure $\al$ on $\Xx$ is Radon if
\(\al(A)=\sup\enscond{\al(K)}{K\subset A\text{ compact}}\)
for every Borel set $A$. We write $\Mm_+(\Xx)$ for finite positive Radon measures and
$\Mm_+^1(\Xx)=\enscond{\al\in\Mm_+(\Xx)}{\al(\Xx)=1}$ for probability measures.
\end{defn}
\(\dotp{f}{\al}\eqdef\int_\Xx f(x)\d\al(x).\)
\(\int_\X f(x) \d\al(x) = \sum_{i=1}^n \a_i f(x_i).\)
\paragraph{Relative densities.}
\begin{defn}[Relative density]\label{def-relative-density}
If $\al$ is absolutely continuous with respect to a positive reference measure $\lambda$, its relative density is the Radon--Nikodym derivative
\(\density{\al}\eqdef\frac{\d\al}{\d\lambda}, \qquad \d\al(x)=\density{\al}(x)\d\lambda(x).\)
Equivalently, for every nonnegative measurable function $h$, and hence for every integrable $h$, \(\int_{\Xx} h(x) \d\al(x) = \int_{\Xx} h(x) \density{\al}(x) \d\lambda(x).\)
\end{defn}
\paragraph{Total variation norm.}
\begin{defn}[Total variation]\label{defn-total-variation}
For a finite signed Radon measure $\al$ on a compact space $\Xx$, \(\norm{\al}_{\TV} \eqdef \usup{f \in \Cc(\Xx)} \enscond{\dotp{f}{\al}}{\norm{f}_\infty \leq 1}.\)
\end{defn}
\(|\al|(A) \eqdef \usup{A=\cup_i B_i} \sum_i |\al(B_i)|,\)
where the supremum is over finite or countable measurable partitions of $A$. Thus, if $\al=\sum_i \a_i \de_{x_i}$ with distinct atoms, $|\al|=\sum_i |\a_i|\de_{x_i}$; if $\d\al(x)=\rho(x)\d\lambda(x)$, then $\d|\al|(x)=|\rho(x)|\d\lambda(x)$.

\begin{prop}[Dual and measure definitions of total variation]\label{prop-tv-dual-measure}
For a finite signed Radon measure $\al$ on a compact space $\Xx$,
one has $\norm{\al}_{\TV}=|\al|(\Xx)$.
\end{prop}
\begin{proof}
The inequality $\norm{\al}_{\TV}\leq |\al|(\Xx)$ follows from
\(\abs{\int f\,\d\al}\leq \int |f|\,\d|\al|\leq \norm{f}_\infty |\al|(\Xx).\)
For the reverse inequality, write the Jordan decomposition $\al=\al^+-\al^-$, so that $|\al|=\al^++\al^-$. The measurable sign $s=\frac{\d\al}{\d|\al|}$ takes values in $\{-1,1\}$ outside a $|\al|$-null set and satisfies $\d\al=s\,\d|\al|$. By regularity of Radon measures on compact spaces, $s$ can be approximated in $L^1(|\al|)$ by continuous functions $f_k$ with $\norm{f_k}_\infty\leq1$. Hence $\int f_k\,\d\al\to\int s\,\d\al=|\al|(\Xx)$, which proves the equality. The final statement is the Riesz--Markov--Kakutani representation theorem with this norm.
\end{proof}
For two absolutely continuous measures $\d\al=\density{\al}\d\lambda$ and $\d\be=\density{\be}\d\lambda$, this gives the concrete formula

\(\norm{\al-\be}_{\TV} = \int_\Xx |\density{\al}(x)-\density{\be}(x)|\,\d\lambda(x).\)
\(\al=\sum_{i=1}^n \tilde a_i\delta_{x_i}, \qquad \be=\sum_{j=1}^m \tilde b_j\delta_{y_j},\)
\(a_k\eqdef\sum_{i:x_i=z_k}\tilde a_i, \qquad b_k\eqdef\sum_{j:y_j=z_k}\tilde b_j.\)
Then $\al=\sum_{k=1}^r a_k\delta_{z_k}$ and $\be=\sum_{k=1}^r b_k\delta_{z_k}$. This convention merges masses located at the same point before comparing the two measures.

\(\norm{\al-\be}_{\TV} = \sum_{k=1}^r |a_k-b_k|.\)
\paragraph{Probabilistic interpretation.}
\begin{defn}[Random variable and law]\label{def-random-variable-law}
Let $(\Omega,\Ff,\PP)$ be a probability space and let $\Xx$ be Polish. An $\Xx$-valued random variable is a measurable map $X:(\Omega,\Ff)\to(\Xx,\Bb(\Xx))$. Its law is
\(\operatorname{Law}(X)\eqdef X_\sharp\PP, \qquad \operatorname{Law}(X)(A)=\PP(X\in A)\)
for every Borel set $A\subset\Xx$.
\end{defn}
For every integrable function $f$, integration with respect to this law is expectation:

\(\int_\X f(x)\,\d\operatorname{Law}(X)(x)=\EE[f(X)].\)
\subsection{Push-Forward}
\label{sec-push-forward}
For a measurable map $\T: \X \rightarrow \Y$, the push-forward operator $\T_\sharp: \Mm(\X) \rightarrow \Mm(\Y)$ records the distribution of the image points.

For discrete measures~\eqref{eq-discr-meas}, the push-forward operation consists simply in moving the positions of all the points in the support of the measure

\(\T_{\sharp} \al \eqdef \sum_i \a_i \de_{\T(x_i)}.\)
\begin{defn}[Push-forward]\label{defn-pushfwd}
For a measurable map $\T: \X \rightarrow \Y$ and $\al \in \Mm(\X)$, the push-forward measure $\be = \T_\sharp \al \in \Mm(\Y)$ is defined, for every measurable set $B\subset\Y$, by
\eql{\label{eq-equiv-pushfwd}
\be(B) = \al\bigl(\enscond{x \in \X}{\T(x) \in B}\bigr).
}
Equivalently, for every bounded measurable function $h:\Y\to\RR$,
\eql{\label{eq-push-fwd}
\int_\Y h(y) \d \be(y) = \int_\X h(\T(x)) \d\al(x).
}
Note that $\T_\sharp$ preserves positivity and total mass, so that if $\al \in \Mm_+^1(\X)$ then $\T_\sharp \al \in \Mm_+^1(\Y)$.
\end{defn}
\begin{prop}[Push-forward formula for densities]\label{prop-push-forward-densities}
Let $U,V\subset\RR^\dims$ be open, let $\al$ and $\be$ have Lebesgue densities $\density{\al}$ and $\density{\be}$ on $U$ and $V$, and let $\T:U\to V$ be a $C^1$ diffeomorphism such that $\be=\T_\sharp\al$. Then, for Lebesgue-almost every $x\in U$,
\eql{\label{eq-pfwd-density}
\density{\al}(x) = |\det(\T'(x))| \density{\be}(\T(x)).
}
Equivalently, for Lebesgue-almost every $y\in V$, \(\density{\be}(y)=\density{\al}(\T^{-1}y)\, \frac{1}{|\det(\T'(\T^{-1}y))|}.\)
\end{prop}
\begin{proof}
Apply~\eqref{eq-push-fwd} to an arbitrary compactly supported continuous test function $h$ and use the change of variables $y=\T(x)$, for which $\d y = |\det(\T'(x))| \d x$.
\begin{align*}
\int_V h(y)\density{\be}(y) \d y &= \int_V h(y) \d \be(y) = \int_U h(\T(x)) \d \al(x) = \int_U h(\T(x)) \density{\al}(x) \d x \\
&= \int_V h(y) \density{\al}(\T^{-1}y) \frac{\d y}{|\det(\T'(\T^{-1} y))|},
\end{align*}
Since the identity holds for every such $h$, uniqueness of Lebesgue densities gives the target-variable formula in the statement. Substituting $y=\T(x)$ then yields~\eqref{eq-pfwd-density}, where $\T'(x)\in\RR^{\dims\times\dims}$ is the Jacobian matrix of $\T$.
\end{proof}
\subsection{Monge's Formulation}
\label{sec-monge-formulation}
\label{sec-continuous-monge}
\paragraph{Monge problem.}
\begin{defn}[Monge problem and Monge map]\label{def-monge-problem}
For $\al\in\Mm_+^1(\Xx)$, $\be\in\Mm_+^1(\Yy)$ and measurable $c:\Xx\times\Yy\to[0,+\infty]$, define
\eql{\label{eq-monge-continuous}
\Monge_c(\al,\be)
\eqdef
\inf_{\T:\Xx\to\Yy\,\text{ measurable}}
\enscond{\int_\Xx c(x,\T(x))\d\al(x)}{\T_\sharp\al=\be}.
}
A feasible $\T$ is a Monge map; a minimizer is an optimal Monge map. The value is $+\infty$ if no feasible map exists.
\end{defn}
\(\T=\cumul{\be}^{-1}\circ\cumul{\al},\)
\begin{prop}[Empirical Monge maps and matchings]\label{prop-empirical-monge-matching}
Assume that the source atoms $x_1,\ldots,x_n$ are distinct and that
\(\al=\frac1n\sum_{i=1}^n\de_{x_i}, \qquad \be=\frac1n\sum_{j=1}^n\de_{y_j}.\)
If $\T_\sharp\al=\be$, then for each distinct target value $z$ in the support of $\be$, exactly $n\be(\{z\})$ source atoms are mapped to $z$. In particular, if the $y_j$ are distinct, then there exists a permutation $\sigma\in\Perm(n)$ such that $\T(x_i)=y_{\sigma(i)}$ for all $i$. Conversely, every such assignment of source atoms to target atoms with the correct masses defines a feasible Monge map on the support of $\al$, and in the distinct-target case
\(\int_\Xx c(x,\T(x))\d\al(x)=\frac1n\sum_{i=1}^n c(x_i,y_{\sigma(i)}).\)
The values of a feasible map away from $\{x_1,\ldots,x_n\}$ may be chosen arbitrarily, provided the resulting extension is measurable. If source locations are repeated, they should first be merged into atoms with larger masses; such atoms cannot be split by a Monge map.
\end{prop}
\begin{proof}
Since $\T_\sharp\al=\be$, all images $\T(x_i)$ must belong to the support of $\be$; otherwise the push-forward would give positive mass to a point outside that support. For any target atom $z$,
\(\be(\{z\})=\al(\T^{-1}(\{z\}))=\frac1n\#\enscond{i}{\T(x_i)=z}.\)
This proves the counting statement. If all target atoms have mass $1/n$, each target receives exactly one source atom, which is a permutation. The converse and the cost identity follow by direct substitution.
\end{proof}
\begin{prop}[Existence of transport maps from atomless sources]\label{prop-existence-transport-map-atomless}
Let $\al$ and $\be$ be Borel probability measures on Polish spaces, and assume that $\al$ is atomless. Then there exists a measurable map $\T$ such that $\T_\sharp\al=\be$.
\end{prop}
\begin{proof}
A standard measure-isomorphism theorem identifies the atomless standard probability space $(\Xx,\al)$ with $([0,1],\mathrm{Leb})$, modulo null sets. It is therefore enough to construct a map from $[0,1]$ to the target law $\be$. Since the target is Polish, it is a standard Borel space; choose a Borel isomorphism $i$ from a full-$\be$ Borel subset of $\Yy$ onto a Borel subset of $[0,1]$. The generalized quantile of $i_\sharp\be$ pushes Lebesgue measure to $i_\sharp\be$. Composing it with $i^{-1}$ and with the source isomorphism gives the desired measurable map, after arbitrary definitions on the discarded null sets.
\end{proof}
\paragraph{Monge distance.}
\begin{defn}[Directed Monge distance]\label{def-directed-monge-distance}
Let $(\Xx,d)$ be a metric space and $1\leq p<+\infty$. For probability measures with finite $p$th moments, define
\eql{\label{eq-monge-distance}
\tilde\Wass_p(\al,\be)^p
\eqdef
\inf_{\T:\Xx\to\Xx\,\text{ measurable}}
\enscond{\int_\Xx d(x,\T(x))^p\d\al(x)}{\T_\sharp\al=\be}.
}
The value is $+\infty$ when the constraint set is empty.
\end{defn}
\begin{prop}[Directed Monge distance]\label{prop-directed-monge-distance}
Let $\Xx=\Yy$ be a Polish metric space and $1\leq p<+\infty$. On probability measures with finite $p$-th moment, the quantity $\tilde\Wass_p$ is nonnegative, vanishes only on the diagonal, and satisfies the triangle inequality. It is therefore a directed extended distance: it need not be symmetric and may take the value $+\infty$.
\end{prop}
\begin{proof}
Nonnegativity is immediate. If $\tilde\Wass_p(\al,\be)=0$, choose feasible maps $\T_k$ with $\int d(x,\T_k(x))^p\d\al(x)\to0$. For every bounded $1$-Lipschitz function $h$,
\[
\left|\int h\d\be-\int h\d\al\right|
=
\left|\int h(\T_k(x))-h(x)\d\al(x)\right|
\leq
\left(\int d(x,\T_k(x))^p\d\al(x)\right)^{1/p}
\longrightarrow0.
\]
Since bounded Lipschitz functions separate probability measures on Polish metric spaces, $\al=\be$. Conversely, the identity map shows that $\tilde\Wass_p(\al,\al)=0$.

If either $\tilde\Wass_p(\al,\ga)$ or $\tilde\Wass_p(\ga,\be)$ is infinite, the claim is automatic. If both are finite, feasible maps $S_\sharp\al=\ga$ and $T_\sharp\ga=\be$ exist, so their composition is feasible from $\al$ to $\be$; in particular, the left-hand side is finite. Fix $\epsilon>0$ and choose $\epsilon$-minimizers satisfying
\(\Ee_\al(S)^{1/p}\leq\tilde\Wass_p(\al,\ga)+\epsilon, \qquad \Ee_\ga(T)^{1/p}\leq\tilde\Wass_p(\ga,\be)+\epsilon.\)
The composed map is feasible from $\al$ to $\be$, and Minkowski's inequality gives
\[
\begin{aligned}
\tilde\Wass_p(\al,\be)
&\leq
\left(\int d(x,T(S(x)))^p\d\al(x)\right)^{1/p} \\
&\leq
\left(\int d(x,S(x))^p\d\al(x)\right)^{1/p}
+
\left(\int d(S(x),T(S(x)))^p\d\al(x)\right)^{1/p} \\
&=
\Ee_\al(S)^{1/p}+\Ee_\ga(T)^{1/p}
\leq
\tilde\Wass_p(\al,\ga)+\tilde\Wass_p(\ga,\be)+2\epsilon.
\end{aligned}
\]
Letting $\epsilon\to0$ gives the result.
\end{proof}
\subsection{Existence and Uniqueness of the Monge Map}
\label{sec-monge-existence-uniqueness}
\paragraph{Brenier's theorem.}
\begin{thm}[Brenier]\label{thm-brenier}
Let $\al,\be\in\Mm_+^1(\RR^d)$ have finite second moments, and assume that $\al$ is absolutely continuous with respect to Lebesgue measure. For the quadratic cost $c(x,y)=\norm{x-y}^2$, there exists a convex function $\phi:\RR^d\to\RR\cup\{+\infty\}$ such that
\(\T=\nabla\phi, \qquad \T_\sharp\al=\be,\)
and $\T$ is the unique optimal Monge map $\al$-almost everywhere. The optimal Kantorovich plan is $(\Id,\T)_\sharp\al$.
\end{thm}
\begin{proof}
Choose optimal $c$-concave potentials $(f,g)$ for the cost $\norm{x-y}^2$ and define
\(\phi(x)\eqdef\frac12\norm{x}^2-\frac12 f(x), \qquad \psi(y)\eqdef\frac12\norm{y}^2-\frac12 g(y).\)
The $c$-concavity of $f$ makes $\phi$ convex. Dual feasibility is equivalent to $\phi(x)+\psi(y)\geq\dotp{x}{y}$, and equality holds for almost every pair under any optimal plan. Replacing $\psi$ by the convex conjugate $\phi^*$ therefore yields the Fenchel equality $\phi(x)+\phi^*(y)=\dotp{x}{y}$ on the optimal relation, hence $y\in\partial\phi(x)$. Since $\al$ has a density and convex functions are differentiable Lebesgue-almost everywhere, $\partial\phi(x)=\{\nabla\phi(x)\}$ for $\al$-almost every $x$. Every optimal plan is therefore concentrated on the same graph, which proves existence and uniqueness of both the optimal plan and the Monge map.
\end{proof}
\(\dotp{\nabla\phi(x)-\nabla\phi(x')}{x-x'}\geq0.\)
\begin{defn}[Measures not charging hypersurfaces]\label{defn-not-charging-hypersurfaces}
A Borel measure $\al$ on $\RR^d$ does not charge hypersurfaces if $\al(S)=0$ for every countably $(d-1)$-rectifiable set $S$, i.e. every set that can be covered, up to an $\mathcal H^{d-1}$-null set, by countably many $C^1$ hypersurfaces.
\end{defn}
\paragraph{Radial measures.}
Radial symmetry gives a useful higher-dimensional case where the Brenier map reduces to a one-dimensional monotone rearrangement. A measure $\al$ on $\RR^d$ is radial if $Q_\sharp\al=\al$ for every orthogonal map $Q$.

\begin{prop}[Optimal transport between radial measures]\label{prop-radial-transport}
Let $\al,\be\in\Mm_+^1(\RR^d)$ be radial probability measures with finite second moments, and assume that $\al$ is absolutely continuous with respect to Lebesgue measure. Define their radial distribution functions
\(F_\al(r)\eqdef \al(\{x:\norm{x}\leq r\}), \qquad F_\be(r)\eqdef \be(\{y:\norm{y}\leq r\}), \qquad r\geq0,\)
and let $F_\be^{-1}(u)\eqdef\inf\{r\geq0:F_\be(r)\geq u\}$ be the generalized inverse.
Set $\tau(r)\eqdef F_\be^{-1}(F_\al(r))$.
Then the quadratic Brenier map from $\al$ to $\be$ is the radial map
\(\T(0)=0, \qquad \T(x)=\frac{\tau(\norm{x})}{\norm{x}}x \quad\text{for }x\neq0.\)
Moreover, if $\al_R=(x\mapsto\norm{x})_\sharp\al$ is the law of the source radius, then
\(\Wass_2^2(\al,\be) = \int_0^\infty |r-\tau(r)|^2\,\d\al_R(r).\)
\end{prop}
\begin{proof}
Let $\al_R=(\norm{\cdot})_\sharp\al$ and $\be_R=(\norm{\cdot})_\sharp\be$ be the radial laws. Since $\al$ is absolutely continuous, $\al_R$ has no atoms. The map $\tau=F_\be^{-1}\circ F_\al$ is therefore the monotone one-dimensional transport from $\al_R$ to $\be_R$. If $X\sim\al$ and $R=\norm{X}$, then $\tau(R)$ has law $\be_R$. Moreover, by radiality, the conditional direction $X/\norm{X}$ given $R=r$ is uniform on the sphere for $\al_R$-a.e. $r>0$. Thus $\T$ changes only the radius, from $R$ to $\tau(R)$, and keeps the uniform angular distribution. It follows that $\T_\sharp\al$ is radial with radial law $\be_R$, hence $\T_\sharp\al=\be$.

The function $\tau$ is nondecreasing and nonnegative. Define $\psi(r)\eqdef\int_0^r\tau(s)\d s$ on $\RR_+$. Then $\psi$ is convex and nondecreasing, so $\phi(x)\eqdef\psi(\norm{x})$ is convex on $\RR^d$. Since $\al$ is absolutely continuous, it does not charge the origin nor the radii where the monotone function $\tau$ is discontinuous. Thus
$\nabla\phi(x)=\tau(\norm{x})x/\norm{x}=\T(x)$ for $\al$-a.e.\ $x$.
The map $\T$ is therefore a gradient of a convex potential and pushes $\al$ to $\be$. Brenier's theorem identifies it as the unique quadratic optimal Monge map. The displayed cost formula follows by substituting $\T(x)$ in $\int\norm{x-\T(x)}^2\d\al(x)$.
\end{proof}
If $\al$ is not absolutely continuous, the same radial idea is still useful but has to be interpreted at the level of couplings of the radial variables: atoms on spheres may have to be split, so a Monge map need not exist.

\paragraph{Polar factorization.}
Brenier's theorem does more than solve one transport problem: it provides a canonical way to extract the ``monotone part'' of a nondegenerate map. Suppose one starts from a square-integrable deformation $u:\Omega\to\RR^d$, for instance a velocity snapshot, an image deformation or a generative map.

\begin{prop}[Polar factorization]\label{prop-polar-factorization}
Let $\Omega\subset\RR^d$ be a bounded domain endowed with normalized Lebesgue measure $\lambda$, and let $u\in L^2(\Omega;\RR^d)$. Assume that the image law $\be=u_\sharp\lambda$ is absolutely continuous with respect to Lebesgue measure. Then there exist a measure-preserving map $s:\Omega\to\Omega$ and a convex function $\phi$ such that
\(u=\nabla\phi\circ s \qquad \lambda\text{-a.e.}\)
The Brenier factor $\nabla\phi$ and the rearrangement $s$ are unique $\lambda$-almost everywhere.
\end{prop}
\begin{proof}
By Brenier's theorem there is a unique convex-gradient map $\T=\nabla\phi$ transporting $\lambda$ to $\be$. Since $\be$ is also absolutely continuous, the reverse Brenier map $S$ transporting $\be$ to $\lambda$ exists, and the two maps are inverse almost everywhere: $\T\circ S=\Id$ $\be$-almost everywhere and $S\circ\T=\Id$ $\lambda$-almost everywhere. Set $s\eqdef S\circ u$. Then
\(s_\sharp\lambda=S_\sharp(u_\sharp\lambda)=S_\sharp\be=\lambda, \qquad \T\circ s=\T\circ S\circ u=u \quad\lambda\text{-a.e.}\)
Thus $s$ is measure preserving and gives the announced factorization. Uniqueness of $\T$ follows from Brenier's theorem, and then $s=S\circ u$ is unique almost everywhere.
\end{proof}
The name is meant to echo the usual polar decomposition of matrices. This analogy becomes literal for linear maps under the Gaussian reference measure.

\paragraph{Displacement interpolation.}
\label{sec-monge-interpolation}
\begin{defn}[Monge and McCann displacement interpolation]\label{def-monge-mccann-interpolation}
If $\T_\sharp\al=\be$, the Monge interpolation generated by $\T$ is the curve
\(\T_t(x)\eqdef(1-t)x+t\T(x), \qquad \al_t\eqdef(\T_t)_\sharp\al, \qquad t\in[0,1].\)
For the quadratic cost, when $\T$ is the optimal Brenier map, this curve is called McCann's displacement interpolation.
\end{defn}
\begin{prop}[Directed Monge displacement geodesics]\label{prop-monge-displacement-geodesic}
Let $1\leq p<+\infty$, let $\al,\be\in\Mm_+^1(\RR^d)$ have finite $p$-th moments, and let $\T$ be an optimal map for $\tilde\Wass_p(\al,\be)$ in~\eqref{eq-monge-distance}. Set $\al_t=(\T_t)_\sharp\al$ with $\T_t=(1-t)\Id+t\T$. Assume that, for every $t<1$, $\T_t$ is one-to-one on a full $\al$-measure Borel set. Its inverse on the image of that set is then measurable and defined $\al_t$-almost everywhere. For $0\leq s\leq t\leq1$,
\(\tilde\Wass_p(\al_s,\al_t)=(t-s)\tilde\Wass_p(\al,\be).\)
Thus $t\mapsto\al_t$ is an oriented constant-speed geodesic for the directed Monge distance. In particular, for $p=2$, this applies to the Brenier map under the hypotheses of Theorem~\ref{thm-brenier}.
\end{prop}
\begin{proof}
The case $s=t$ is trivial, so assume $s<t$. Since $s<1$, the inverse $\T_s^{-1}$ is defined $\al_s$-almost everywhere along the transported particles. Define
$S_{s,t}\eqdef\T_t\circ\T_s^{-1}$ $\al_s$-a.e.
Then $(S_{s,t})_\sharp\al_s=\al_t$, and, using the optimality of $\T$, \(\tilde\Wass_p(\al_s,\al_t)^p \leq \int \norm{S_{s,t}(z)-z}^p\d\al_s(z) = \int \norm{\T_t(x)-\T_s(x)}^p\d\al(x) = (t-s)^p\tilde\Wass_p(\al,\be)^p.\)
The same particle construction gives $\tilde\Wass_p(\al,\al_s)\leq s\tilde\Wass_p(\al,\be)$. If $t<1$, the map $\T\circ\T_t^{-1}$ sends $\al_t$ to $\be$ and gives $\tilde\Wass_p(\al_t,\be)\leq(1-t)\tilde\Wass_p(\al,\be)$; if $t=1$, this latter distance is zero. Using the triangle inequality from Proposition~\ref{prop-directed-monge-distance},
\[
\tilde\Wass_p(\al,\be)
\leq
\tilde\Wass_p(\al,\al_s)+\tilde\Wass_p(\al_s,\al_t)+\tilde\Wass_p(\al_t,\be)
\leq
s\tilde\Wass_p(\al,\be)+\tilde\Wass_p(\al_s,\al_t)+(1-t)\tilde\Wass_p(\al,\be).
\]
Hence $\tilde\Wass_p(\al_s,\al_t)\geq(t-s)\tilde\Wass_p(\al,\be)$, which proves equality. For a Brenier map $\T=\nabla\phi$, the map $\T_t$ is the gradient of
\(x\mapsto (1-t)\frac{\norm{x}^2}{2}+t\phi(x)\),
which is $(1-t)$-strongly convex for every $t<1$. On the full-measure set where $\phi$ is differentiable, this gives
$\dotp{\T_t(x)-\T_t(y)}{x-y}\geq(1-t)\norm{x-y}^2$, so $\T_t$ is injective there.
\end{proof}
\paragraph{Regularity and the Monge--Amp\`ere equation.}
\begin{prop}[Caffarelli regularity]\label{prop-caffarelli-regularity}
Let $\Omega,\Lambda\subset\RR^d$ be bounded uniformly convex domains with $C^2$ boundaries. Let $\al=\density{\al}(x)\d x$ and $\be=\density{\be}(y)\d y$ be supported on $\Omega$ and $\Lambda$, respectively, with
\(0<m\leq\density{\al},\density{\be}\leq M<+\infty.\)
If $\density{\al}\in C^\alpha(\overline\Omega)$ and $\density{\be}\in C^\alpha(\overline\Lambda)$ for some $\alpha\in(0,1)$, then the Brenier potential $\phi$ is strictly convex in $\Omega$ and belongs to $C^{2,\alpha}_{\mathrm{loc}}(\Omega)$. Under the standard compatibility and higher boundary-regularity assumptions, the same regularity extends to the boundary.
\end{prop}
\begin{proof}
The potential solves the Monge--Amp\`ere equation
\(\det(\nabla^2\phi(x)) =\frac{\density{\al}(x)}{\density{\be}(\nabla\phi(x))}\)
in the Alexandrov sense, with second boundary condition $\nabla\phi(\Omega)=\Lambda$. The density bounds and convexity of the domains give strict convexity and localization of sections. Caffarelli's interior theory then yields $C^{2,\alpha}_{\mathrm{loc}}$ estimates for $\phi$; the boundary statement follows from the boundary regularity theory under uniform convexity and compatibility assumptions.
\end{proof}
For sufficiently smooth positive densities and a classical Brenier map, the change-of-variables formula~\eqref{eq-pfwd-density} gives the Monge--Amp\`ere equation

\eql{\label{eq-monge-ampere}
\det(\nabla^2\phi(x))\density{\be}(\nabla\phi(x))=\density{\al}(x).
}
\begin{prop}[Linearization of the Monge--Amp\`ere equation]\label{prop-linearized-monge-ampere}
Let $\rho_\epsilon=\rho_0+\epsilon r+o(\epsilon)$ be a smooth perturbation of a positive reference density $\rho_0$ on a smooth bounded domain $\Omega$, with $\int_\Omega r\d x=0$. Assume that the expansions hold in a topology strong enough to differentiate the change-of-variables identity. If $\T_\epsilon(x)=x+\epsilon\nabla u(x)+o(\epsilon)$ transports $\rho_0\d x$ to $\rho_\epsilon\d x$, then, to first order,
\(-\nabla\cdot(\rho_0\nabla u)=r.\)
In particular, when $\rho_0$ is constant, the linearized equation is $-\Delta u=r/\rho_0$. If the source and target occupy the same fixed domain and no mass crosses its boundary, the accompanying first-order boundary condition is $\rho_0\partial_n u=0$ on $\partial\Omega$; together with $\int_\Omega u\d x=0$, it fixes the additive constant.
\end{prop}
\begin{proof}
The change-of-variables equation for $\T_\epsilon$ is \(\rho_0(x)=\rho_\epsilon(\T_\epsilon(x))\det(\nabla \T_\epsilon(x)).\)
Expanding $\rho_\epsilon(x+\epsilon\nabla u)=\rho_0(x)+\epsilon r(x)+\epsilon\dotp{\nabla\rho_0}{\nabla u}+o(\epsilon)$ and
$\det(\Id+\epsilon\nabla^2u)=1+\epsilon\Delta u+o(\epsilon)$ gives
\[
\rho_0
=
\rho_0+
\epsilon\bigl(r+\dotp{\nabla\rho_0}{\nabla u}+\rho_0\Delta u\bigr)+o(\epsilon)
=
\rho_0+
\epsilon\bigl(r+\nabla\cdot(\rho_0\nabla u)\bigr)+o(\epsilon).
\]
The first-order term must vanish.
\end{proof}
\subsection{Beyond the Quadratic Euclidean Cost}
\label{sec-beyond-quadratic-euclidean-cost}
\paragraph{\texorpdfstring{$\Wass_p$ costs.}{Wp costs}.}
The quadratic cost is special because it identifies the optimal map with the Euclidean gradient of an ordinary convex potential. For the Wasserstein cost, normalized as \(c(x,y)=\norm{x-y}^p/p\) with \(p>1\), the same Monge-map picture survives, but the potential is adapted to the cost.

\(\nabla f(x)=\nabla_x c(x,\T(x)).\)
\(\T(x)=x-\norm{\nabla f(x)}^{q-2}\nabla f(x), \qquad \frac1p+\frac1q=1.\)
\paragraph{Squared geodesic distance.}
\(\T(x)=\exp_x(-\nabla\phi(x)),\)
\[
\T_t(x)
=
\exp_x\!\left(t\exp_x^{-1}(\T(x))\right)
=
\exp_x(-t\nabla\phi(x)),
\qquad
\al_t=(\T_t)_\sharp\al,
\qquad t\in[0,1].
\]
\paragraph{Poincar\'e disk.}
\paragraph{Twist condition.}
\begin{defn}[Twist condition]\label{def-twist-condition}
Let $X,Y\subset\RR^d$ and let $c\in C^1(X\times Y)$. The cost $c$ satisfies the twist condition from $X$ to $Y$ if, for every fixed $x\in X$, the map
\(y\in Y \longmapsto \nabla_x c(x,y)\in\RR^d\)
is injective. The dual twist condition is the analogous injectivity of $x\mapsto\nabla_y c(x,y)$ for every fixed $y$.
\end{defn}
\begin{prop}[Twist prevents splitting]\label{prop-twist-prevents-splitting}
Assume that $c\in C^1(X\times Y)$ satisfies the twist condition from $X$ to $Y$. Suppose that optimal pairs for the cost $c$ admit a source potential $f$ with the following first-order certificate: $f$ is differentiable on a full $\al$-measure subset of $X$, and whenever $x$ belongs to this differentiability set and $y$ can be paired optimally with $x$, one has $\nabla f(x)=\nabla_x c(x,y)$. Then, for $\al$-almost every source point $x$, there is at most one target point $y$ that can be paired optimally with $x$.
\end{prop}
\begin{proof}
Let $(f,g)$ be optimal Kantorovich potentials and let $\pi$ be an optimal plan. Complementary slackness gives $f(x)+g(y)=c(x,y)$ for $\pi$-almost every $(x,y)$, while dual feasibility gives $f(z)+g(y)\leq c(z,y)$ for all $z$. Thus, for almost every contact pair $(x,y)$, the function $z\mapsto c(z,y)-g(y)$ touches $f$ from above at $x$. At a point where $f$ is differentiable, this gives
\(\nabla f(x)=\nabla_x c(x,y)\).
If the cost is twisted, this equation determines at most one $y$. Hence no optimal plan can split the mass of such an $x$ between several target points. Since differentiability holds on a full $\al$-measure set, the plan is concentrated on the graph of a measurable map there.
\end{proof}
The quadratic cost satisfies twist since $\nabla_x\norm{x-y}^2=2(x-y)$; the bilinear cost $c(x,y)=-\dotp{x}{y}$ satisfies twist since $\nabla_x c(x,y)=-y$; more generally $c(x,y)=h(x-y)$ is twisted when $h$ is smooth strictly convex and $\nabla h$ is injective. On a Riemannian manifold, the squared geodesic cost is twisted locally away from the cut locus.

\paragraph{Ma--Trudinger--Wang curvature.}
\begin{defn}[Weak MTW condition]\label{def-mtw-condition}
Assume that $c\in C^4(X\times Y)$, that $X,Y\subset\RR^d$, that $c$ is twisted, and that the mixed Hessian $\nabla^2_{xy}c(x,y)$ is invertible. For $(x,p)$ in the image of the change of variables $(x,y)\mapsto(x,-\nabla_x c(x,y))$, let $Y(x,p)$ be defined by
$-\nabla_x c(x,Y(x,p))=p$, and set
$A_{i,j}(x,p)\eqdef-\partial^2_{x_i x_j}c(x,Y(x,p))$.
The cost satisfies the weak MTW condition if
\(\sum_{i,j,k,l}\partial^2_{p_k p_l} A_{i,j}(x,p)\,\xi_i\xi_j\eta_k\eta_l\geq0 \qquad\text{whenever }\dotp{\xi}{\eta}=0.\)
A strong MTW condition asks for a positive lower bound, proportional to $\norm{\xi}^2\norm{\eta}^2$, on the same orthogonal directions.
\end{defn}
\begin{prop}[MTW condition and regularity]\label{prop-mtw-regularity-implication}
Assume that $c$ is smooth, twisted and non-degenerate, that the source and target domains satisfy the appropriate mutual $c$-convexity assumptions, and that the source and target densities are positive and bounded above and below. Under the weak MTW condition, the optimal map is locally H\"older continuous; with the corresponding global domain hypotheses, it is continuous up to the boundary. Under the strong MTW condition and higher smoothness of the data and domains, one obtains higher regularity estimates.
\end{prop}
\begin{proof}
This is the regularity theory of Ma--Trudinger--Wang, Trudinger--Wang and Loeper, stated in structural form rather than with its full boundary hypotheses. The $c$-convex potential solves a generated Jacobian equation. Twist and non-degeneracy turn its optimal relation into a single-valued map, while the weak MTW inequality controls the geometry of contact sets and yields interior H\"older estimates. Strong MTW provides a quantitative curvature lower bound; combined with smooth densities, domain geometry and elliptic estimates, it yields higher regularity. Conversely, Loeper proved that weak MTW is necessary for continuity for arbitrary smooth positive data.
\end{proof}
\subsection{One-Dimensional Transport and Quantiles}
\label{sec-1d-transport-quantiles}
\paragraph{Cumulative and quantile functions.}
\begin{defn}[Cumulative and quantile functions]\label{def-cdf-quantile}
For $\al\in\Mm_+^1(\RR)$, its cumulative distribution function is
\eql{\label{eq-cumul-defn}
\cumul{\al}(x)\eqdef\al((-\infty,x]).
}
Its generalized inverse, or quantile function, is
\eql{\label{eq-OT-map-1d}
\cumul{\al}^{-1}(r)
\eqdef
\inf\enscond{x\in\RR}{\cumul{\al}(x)\geq r},
\qquad r\in(0,1).
}
\end{defn}
\begin{prop}[Quantile push-forward]\label{prop-quantile-pushforward}
One has $(\cumul{\al}^{-1})_\sharp\mathrm{Leb}_{[0,1]}=\al$. If $\al$ has no atoms, then $(\cumul{\al})_\sharp\al=\mathrm{Leb}_{[0,1]}$.
\end{prop}
\begin{proof}
For simplicity, assume first that $\al$ has a strictly positive density, so that $\cumul{\al}$ is strictly increasing and continuous. Denote $\ga\eqdef(\cumul{\al}^{-1})_\sharp\mathrm{Leb}_{[0,1]}$. It is enough to prove that $\cumul{\ga}=\cumul{\al}$. For every $x$,
\[
\begin{aligned}
\cumul{\ga}(x)
&=
\int_0^1 \mathbf 1_{(-\infty,x]}(\cumul{\al}^{-1}(z))\d z
=
\int_0^1 \mathbf 1_{[0,\cumul{\al}(x)]}(z)\d z
=
\cumul{\al}(x),
\end{aligned}
\]
where we used $\cumul{\al}^{-1}(z)\leq x$ if and only if $z\leq\cumul{\al}(x)$. General measures follow from the same argument with generalized inverses and right-continuity of the cumulative distribution function. If $\al$ has no atoms, the probability integral transform gives $(\cumul{\al})_\sharp\al=\mathrm{Leb}_{[0,1]}$.
\end{proof}
\paragraph{1-D Monge solutions.}\label{par-1d-monge-solution}
\begin{thm}[One-dimensional Monge solution]\label{prop-1d-quantile-map}
Let $\al,\be\in\Mm_+^1(\RR)$, assume that $\al$ has no atoms, and let $h:\RR\to\RR$ be convex. Consider the cost $c(x,y)=h(x-y)$. Write $q_\al\eqdef\cumul{\al}^{-1}$ and $q_\be\eqdef\cumul{\be}^{-1}$, and assume that $h(q_\al-q_\be)\in L^1(0,1)$. Then
\eql{\label{eq-1d-monge-map}
\T=q_\be\circ\cumul{\al}
=
\cumul{\be}^{-1}\circ\cumul{\al}
}
satisfies $\T_\sharp\al=\be$ and minimizes
\(\umin{S_\sharp\al=\be}\int h(x-S(x))\d\al(x).\)
In particular, $h(t)=|t|^p$ gives the usual one-dimensional Monge map for every $1\leq p<+\infty$ whenever the source has no atoms.
\end{thm}
\begin{proof}
By Proposition~\ref{prop-quantile-pushforward}, atomlessness of $\al$ gives $(\cumul{\al})_\sharp\al=\mathrm{Leb}_{[0,1]}$, while $(q_\be)_\sharp\mathrm{Leb}_{[0,1]}=\be$. Hence $\T_\sharp\al=\be$.

Let $S$ be any admissible map and set $\pi_S=(\Id,S)_\sharp\al$. The one-dimensional uncrossing theorem for relaxed couplings, proved later as Theorem~\ref{prop-1d-kantorovich-quantile-coupling}, shows that the quantile coupling $\pi^\star=(q_\al,q_\be)_\sharp\mathrm{Leb}_{[0,1]}$ is optimal among all couplings. Since $\al$ has no atoms, $(\Id,\T)_\sharp\al=\pi^\star$, and therefore
\(\int h(x-\T(x))\d\al(x) = \int h(x-y)\d\pi^\star(x,y) \leq \int h(x-y)\d\pi_S(x,y) = \int h(x-S(x))\d\al(x).\)
This proves optimality of $\T$ among Monge maps.
\end{proof}
For discrete measures, one cannot directly apply the map formula when the source has atoms, but if the measures are uniform on the same number of Dirac masses, then it is exactly the sorting formula of Proposition~\ref{prop-matching-1d-monotone}.

\begin{prop}[One-dimensional Wasserstein formulas]\label{prop-wass-quantile-1d}
Let $1\leq p<+\infty$ and let $\al,\be\in\Mm_+^1(\RR)$ have finite $p$-th moments. Then
\eql{\label{eq-wass-cumul}
\Wass_p(\al,\be)^p
=
\int_0^1\abs{\cumul{\al}^{-1}(r)-\cumul{\be}^{-1}(r)}^p\d r
=
\norm{\cumul{\al}^{-1}-\cumul{\be}^{-1}}_{L^p([0,1])}^p.
}
For $p=1$, this is equivalently
\begin{equation}\label{eq-w1-1d}
\Wass_1(\al,\be)
=
\norm{\cumul{\al}-\cumul{\be}}_{L^1(\RR)}
=
\int_\RR \abs{\cumul{\al}(x)-\cumul{\be}(x)}\d x
=
\int_\RR \abs{\int_{-\infty}^x\d(\al-\be)}\d x.
\end{equation}
\end{prop}
\begin{proof}
The first formula follows from Theorem~\ref{prop-1d-kantorovich-quantile-coupling}: the optimal coupling is obtained by taking the same quantile level $r$ for both measures. For $p=1$, use the layer-cake identity. If $q_\al$ and $q_\be$ are the two quantile functions, then
\(\int_0^1 |q_\al(r)-q_\be(r)|\d r = \int_\RR \lambda\bigl(\{r:q_\al(r)\leq x<q_\be(r)\}\cup\{r:q_\be(r)\leq x<q_\al(r)\}\bigr)\d x.\)
The measure of the displayed set is exactly $|\cumul{\al}(x)-\cumul{\be}(x)|$ for almost every $x$.
\end{proof}
\begin{prop}[Bobkov--Ledoux cumulative formula]\label{prop-bobkov-ledoux-cdf-w2}
Let $\al,\be\in\Mm_+^1([a,b])$ and write $s_+\eqdef\max(s,0)$. Then
\eql{\label{eq-bobkov-ledoux-cdf-w2}
\Wass_2^2(\al,\be)
=
2\int_a^b\int_x^b
\Big[
\big(\cumul{\al}(x)-\cumul{\be}(y)\big)_+
+
\big(\cumul{\be}(x)-\cumul{\al}(y)\big)_+
\Big]\d y\,\d x.
}
\end{prop}
\begin{proof}
For any $u,v\in[a,b]$, one has \(|u-v|^2 = 2\int_a^b\int_x^b \Big[ \mathbf 1_{\{u\leq x\leq y<v\}} + \mathbf 1_{\{v\leq x\leq y<u\}} \Big]\d y\,\d x,\)
because, when $u<v$, the integration domain is the triangle $u\leq x\leq y<v$, whose area is $(v-u)^2/2$, and the case $v<u$ is symmetric. Apply this identity with $u=q_\al(r)$ and $v=q_\be(r)$, where $q_\al=\cumul{\al}^{-1}$ and $q_\be=\cumul{\be}^{-1}$, and integrate with respect to $r\in(0,1)$. Fubini's theorem applies since the integrand is bounded. For $x\leq y$, the generalized inverse identities give, up to endpoint sets of zero Lebesgue measure,
\(\lambda\enscond{r}{q_\al(r)\leq x\ \text{and}\ q_\be(r)>y} = \big(\cumul{\al}(x)-\cumul{\be}(y)\big)_+,\)
and the same argument with $\al$ and $\be$ exchanged gives the second positive part.
\end{proof}
\(\cumul{\al_t}^{-1}(r) = (1-t)\cumul{\al}^{-1}(r)+t\cumul{\be}^{-1}(r), \qquad r\in(0,1).\)
\paragraph{OT on trees.}
\begin{prop}[Cumulative formula on a tree]\label{prop-tree-w1-cumulative}
Let $\mathsf{T}=(V,E,\ell)$ be a finite weighted tree, where each edge $e$ has length $\ell_e>0$, and let $d_{\mathsf{T}}$ be its geodesic distance. Choose a root $o$. For an edge $e=(u,v)$ oriented away from $o$, let $V_e$ be the set of vertices in the connected component containing $v$ after removing $e$. If $\al,\be\in\Pp(V)$, then
\eql{\label{eq-tree-w1-cumulative}
\Wass_{1,d_{\mathsf{T}}}(\al,\be)
=
\umin{\pi\in\Couplings(\al,\be)}
\sum_{x,y\in V} d_{\mathsf{T}}(x,y)\pi_{x y}
=
\sum_{e\in E}\ell_e
\abs{\al(V_e)-\be(V_e)}.
}
The value is computable in $O(|V|)$ arithmetic operations once the signed masses $\al-\be$ are stored on the vertices, by one postorder traversal of the rooted tree.
\end{prop}
\begin{proof}
Let $\sigma=\al-\be$. Removing an edge $e$ splits the tree into $V_e$ and its complement. For any coupling $\pi$, the net amount of mass crossing $e$ from $V_e$ to $V\setminus V_e$ minus the amount crossing in the reverse direction is fixed and equals $\sigma(V_e)$. Hence the total amount of transported mass whose path uses $e$ is at least $|\sigma(V_e)|$. Since every unit of mass crossing $e$ pays the length $\ell_e$, summing over edges gives
\(\sum_{x,y}d_{\mathsf{T}}(x,y)\pi_{x,y} = \sum_{e\in E}\ell_e \sum_{x,y:\ e\in[x,y]}\pi_{x,y} \geq \sum_{e\in E}\ell_e|\sigma(V_e)|,\)
where $[x,y]$ is the unique path between $x$ and $y$.

Process the rooted tree from the leaves to the root. At a vertex $v\neq o$, compute the total signed mass in the subtree below $v$, namely $F_e=\sigma(V_e)$ for the edge $e$ joining $v$ to its parent. If $F_e>0$, send $F_e$ units of surplus from this subtree to the parent side; if $F_e<0$, import $-F_e$ units from the parent side into the subtree. After this operation the subtree has zero net imbalance and can be collapsed into its parent. Continuing upward balances all subtrees because $\sigma(V)=0$. Decomposing these edge fluxes into source-to-target paths yields a coupling whose traffic across each edge is exactly $|F_e|$, and therefore whose cost is the right-hand side of~\eqref{eq-tree-w1-cumulative}. The same postorder pass computes all $F_e$, proving the complexity claim.
\end{proof}
\paragraph{Triangular rearrangements.}
\begin{prop}[Knothe--Rosenblatt triangular rearrangement]\label{prop-knothe-rosenblatt}
Let $\al,\be\in\Mm_+^1(\RR^d)$. Assume that the first marginal of $\al$ and, recursively, the one-dimensional conditional source laws used below are atomless almost everywhere. Fix measurable versions of all regular conditional distributions. Then there is a triangular map
\(\T(x_1,\ldots,x_d) = (\T_1(x_1),\T_2(x_1,x_2),\ldots,\T_d(x_1,\ldots,x_d))\)
such that $\T_\sharp\al=\be$ and, for each $k$, the function $x_k\mapsto\T_k(x_1,\ldots,x_k)$ is nondecreasing for $\al_{<k}$-almost every $x_{<k}=(x_1,\ldots,x_{k-1})$.
\end{prop}
\begin{proof}
The construction is recursive. For $k=1$, let $\T_1$ be the monotone rearrangement between the first marginal of $\al$ and the first marginal of $\be$. Suppose that $\T_1,\ldots,\T_{k-1}$ have already been constructed. Write $x_{<k}=(x_1,\ldots,x_{k-1})$ and $\T_{<k}=(\T_1,\ldots,\T_{k-1})$. By induction, $(\T_{<k})_\sharp\al_{<k}=\be_{<k}$, where $\al_{<k}$ and $\be_{<k}$ are the first $(k-1)$-coordinate marginals. Let $\al^k_{x_{<k}}$ and $\be^k_{y_{<k}}$ be regular conditional laws of the $k$-th coordinate given the previous coordinates. Define $\T_k(x_{<k},\cdot)$ as the one-dimensional monotone rearrangement from $\al^k_{x_{<k}}$ to $\be^k_{\T_{<k}(x_{<k})}$.

The map $\T_k(x_{<k},\cdot)$ is nondecreasing by the one-dimensional rearrangement theorem. The chain rule for disintegrations then shows that, after step $k$, the first $k$ coordinates of $\T_\sharp\al$ have the same law as the first $k$ coordinates of $\be$. At $k=d$ this gives $\T_\sharp\al=\be$. Target atoms are handled by generalized quantiles. If a source conditional law has atoms that must be split, the same recursive construction defines a triangular Markov kernel rather than a deterministic map.
\end{proof}
\begin{prop}[Anisotropic Brenier maps converge to Knothe--Rosenblatt]\label{prop-knothe-limit-anisotropic-brenier}
Let $\al,\be$ be compactly supported probability measures on $\RR^d$ with densities bounded above and below on their rectangular supports. In particular, the source and target conditional laws entering the Knothe--Rosenblatt construction are atomless almost everywhere. For $\epsilon>0$, set
\(c_\epsilon(x,y) \eqdef \sum_{k=1}^d \epsilon^{k-1}\abs{x_k-y_k}^2.\)
Let $T_\epsilon$ be the Monge map from $\al$ to $\be$ for the cost $c_\epsilon$, and let $T_{\mathrm{KR}}$ be the triangular Knothe--Rosenblatt rearrangement with the coordinate order used above. Then
\(T_\epsilon \longrightarrow T_{\mathrm{KR}} \qquad\text{in } L^2(\al;\RR^d) \qquad\text{as } \epsilon\to0.\)
\end{prop}
\begin{proof}
We summarize the scale-separated argument. Let $\pi_\epsilon=(\Id,T_\epsilon)_\sharp\al$. For a coupling $\gamma$, set $I_k(\gamma)\eqdef\int\abs{x_k-y_k}^2\d\gamma(x,y)$ and $F_\epsilon(\gamma)\eqdef\sum_{k=1}^d\epsilon^{k-1}I_k(\gamma)$.
Since the supports are compact, every sequence $\epsilon\downarrow0$ has a subsequence for which $\pi_\epsilon$ converges weakly to a coupling $\pi$ between $\al$ and $\be$. The coordinate costs are bounded and continuous on the common compact support. Dividing $F_\epsilon$ by its leading coefficient and passing to the limit shows that $\pi$ minimizes $I_1$. The first source marginal is atomless, so the $(x_1,y_1)$-marginal of $\pi$ is the unique monotone coupling; equivalently, $y_1=T_1(x_1)$ under $\pi$.

Optimality against the Knothe coupling $\pi_{\mathrm{KR}}$ gives $F_\epsilon(\pi_\epsilon)\leq F_\epsilon(\pi_{\mathrm{KR}})$. Since $I_1(\pi_\epsilon)$ is bounded below by the optimal first-coordinate value $I_1(\pi_{\mathrm{KR}})$, cancellation and division by $\epsilon$ yield
\(\sum_{k=2}^d\epsilon^{k-2}I_k(\pi_\epsilon) \leq\sum_{k=2}^d\epsilon^{k-2}I_k(\pi_{\mathrm{KR}}).\)
Passing to the limit shows that $\pi$ minimizes $I_2$ among couplings with the already identified first-coordinate marginal. Disintegration with respect to $(x_1,y_1)$ turns this into the monotone transport between the conditional laws of $x_2$ and $y_2$, so $y_2=T_2(x_1,x_2)$ under $\pi$.

The higher-coordinate induction repeats this comparison after disintegrating over the already identified coordinates. Its technical point is to control the residual suboptimality of the earlier coordinates before dividing by the next power of $\epsilon$; the atomlessness of both source and target conditionals makes the preceding monotone couplings invertible almost everywhere. The scale-separated estimate proved in then gives successively $y_k=T_k(x_1,\ldots,x_k)$ for every $k$.

Finally, let $X\sim\al$. The graph couplings are the laws of $(X,T_\epsilon(X))$, and they converge weakly to the law of $(X,T_{\mathrm{KR}}(X))$. To turn this into convergence of maps, fix $\zeta>0$. By Lusin's theorem, there is a compact set $K$ with $\al(K)>1-\zeta$ on which $T_{\mathrm{KR}}$ is continuous. On $K$, the set
\(\enscond{(x,y)}{x\in K,\ \norm{y-T_{\mathrm{KR}}(x)}\geq\delta}\)
is closed and has zero mass under the limiting graph coupling. Portmanteau's theorem gives
\[
\limsup_{\epsilon\to0}
\al\!\left(\enscond{x}{x\in K,\ \norm{T_\epsilon(x)-T_{\mathrm{KR}}(x)}\geq\delta}\right)
=0.
\]
Adding the complement of $K$ and letting $\zeta\to0$ proves convergence in $\al$-probability. Since all maps take values in a common compact set, their squared differences are uniformly integrable, and convergence therefore holds in $L^2(\al)$.
\end{proof}
\subsection{Gaussian Measures and the Bures Metric}
\label{sec-gaussian-bures}
\paragraph{One-dimensional Gaussians.}
Let $\al=\Gaussian(m_\al,\sigma_\al^2)$ and $\be=\Gaussian(m_\be,\sigma_\be^2)$ be two nondegenerate Gaussians on $\RR$. Then

\(\T(x)=\frac{\sigma_\be}{\sigma_\al}(x-m_\al)+m_\be\)
satisfies $\T_\sharp\al=\be$. It is the derivative of the convex function

\(\phi(x)=\frac{\sigma_\be}{2\sigma_\al}(x-m_\al)^2+m_\be x,\)
\[
\Wass_2(\al,\be)^2
=
\int_\RR\left(\frac{\sigma_\be}{\sigma_\al}(x-m_\al)+m_\be-x\right)^2\d\al(x)
=
(m_\al-m_\be)^2+(\sigma_\al-\sigma_\be)^2.
\]
The formula extends by continuity to Dirac masses, although the affine Monge map itself only pushes a Dirac source to another Dirac. Thus the OT geometry of one-dimensional Gaussians is the Euclidean geometry of the closed half-plane $(m,\sigma)\in\RR\times\RR_+$.

\paragraph{Multivariate Gaussians.}
\eql{\label{eq-gauss-pf}
\al=\Gaussian(\mean_\al,\cov_\al),
\qquad
\be=\Gaussian(\mean_\be,\cov_\be),
\qquad
\T(x)=\mean_\be+A(x-\mean_\al),
}
\begin{prop}[Affine push-forward of Gaussians]\label{prop-gaussian-affine-push-forward}
One has $\T_\sharp\al=\be$ if and only if
\eql{\label{eq-gauss-covariance-pushforward}
A\cov_\al A^\top=\cov_\be.
}
\end{prop}
\begin{proof}
An affine function maps a Gaussian to a Gaussian, so the law of $\T(X)$ is determined by its mean and covariance. If $X\sim\al$ and $Y=\T(X)$, then
\(\EE(Y)=\mean_\be+A\EE(X-\mean_\al)=\mean_\be\),
and
\(\EE((Y-\mean_\be)(Y-\mean_\be)^\top) = A\EE((X-\mean_\al)(X-\mean_\al)^\top)A^\top = A\cov_\al A^\top.\)
Thus $A\cov_\al A^\top=\cov_\be$ is necessary and sufficient for $Y$ to have the same mean and covariance as $\be$; since both laws are Gaussian, this is equivalent to equality in distribution.
\end{proof}
\begin{defn}[Bures metric]\label{def-bures-metric}
For positive semidefinite covariance matrices $\Sigma$ and $\Lambda$, the Bures metric is
\eql{\label{eq-bure-defn}
\Bb(\Sigma,\Lambda)^2
\eqdef
\tr\pa{\Sigma+\Lambda-2(\Sigma^{1/2}\Lambda\Sigma^{1/2})^{1/2}}.
}
\end{defn}
\begin{prop}[Gaussian $\Wass_2$ formula and Bures covariance term]\label{prop-gaussian-w2-bures}
Assume that $\cov_\al$ and $\cov_\be$ are positive definite. The covariance equation and its unique symmetric positive-definite solution are
\eql{\label{eq-bures-map}
A\cov_\al A=\cov_\be,\quad A\succ0,
\qquad\Longleftrightarrow\qquad
A=
\cov_\al^{-1/2}
\pa{\cov_\al^{1/2}\cov_\be\cov_\al^{1/2}}^{1/2}
\cov_\al^{-1/2}.
}
The affine map $\T(x)=\mean_\be+A(x-\mean_\al)$ is the optimal quadratic-cost transport from $\Gaussian(\mean_\al,\cov_\al)$ to $\Gaussian(\mean_\be,\cov_\be)$, and
\eql{\label{eq-dist-gauss}
\Wass_2(\al,\be)^2
=
\norm{\mean_\al-\mean_\be}^2+\Bb(\cov_\al,\cov_\be)^2,
}
where $\Bb$ is the Bures metric of Definition~\ref{def-bures-metric}.
\end{prop}
\begin{proof}
Multiplying $A\cov_\al A=\cov_\be$ on the left and right by $\cov_\al^{1/2}$ gives
$(\cov_\al^{1/2}A\cov_\al^{1/2})^2
=\cov_\al^{1/2}\cov_\be\cov_\al^{1/2}$.
Since $A$ is symmetric positive, $\cov_\al^{1/2}A\cov_\al^{1/2}$ is symmetric positive and is therefore the unique positive square root of the right-hand side. Conversely, the matrix in~\eqref{eq-bures-map} is symmetric positive and satisfies the covariance equation.

By Proposition~\ref{prop-gaussian-affine-push-forward}, this affine map pushes $\al$ to $\be$. It is the gradient of a convex quadratic potential, so Brenier's theorem implies optimality. If $X\sim\al$, then
\begin{align*}
\EE\norm{X-\T(X)}^2
&=
\norm{\mean_\al-\mean_\be}^2
+
\EE\norm{(\Id-A)(X-\mean_\al)}^2 \\
&=
\norm{\mean_\al-\mean_\be}^2
+
\tr((\Id-A)\cov_\al(\Id-A)^\top) \\
&=
\norm{\mean_\al-\mean_\be}^2
+
\tr(\cov_\al)+\tr(A\cov_\al A)-2\tr(A\cov_\al) \\
&=
\norm{\mean_\al-\mean_\be}^2
+
\tr(\cov_\al)+\tr(\cov_\be)
-2\tr((\cov_\al^{1/2}\cov_\be\cov_\al^{1/2})^{1/2}),
\end{align*}
which is the desired expression. The formula for singular covariance matrices follows by adding $\eta\Id$ and letting $\eta\downarrow0$.
\end{proof}
\paragraph{Alternate formulation of the Bures metric.}
\begin{prop}[Bures metric as a second-moment quotient]\label{prop-bures-second-moment-lift}
For $\al\in\Pp_2(\RR^d)$, set its raw second-moment matrix to be
\(\Phi_2(\al)\eqdef \int_{\RR^d} uu^\top \d\al(u)\in\mathbb S_+^d.\)
For $\Sigma,\Lambda\in\mathbb S_+^d$, \(\min_{\Phi_2(\al)=\Sigma,\ \Phi_2(\be)=\Lambda} \Wass_2(\al,\be)^2 = \Bb(\Sigma,\Lambda)^2.\)
In other words, the Bures metric is the Wasserstein quotient distance induced on positive semidefinite matrices by identifying all probability measures with the same raw second-moment matrix.
\end{prop}
\begin{proof}
Assume first that $\Sigma$ and $\Lambda$ are positive definite. Fix two admissible laws $\al,\be$ and an arbitrary coupling $\pi$ between them, and write
$K\eqdef\int_{\RR^d\times\RR^d}uv^\top\d\pi(u,v)$. Then
$\int\norm{u-v}^2\d\pi(u,v)
=\tr(\Sigma)+\tr(\Lambda)-2\tr(K)$.
The block second-moment matrix
\[
\int
\begin{pmatrix}u\\ v\end{pmatrix}
\begin{pmatrix}u\\ v\end{pmatrix}^{\!\top}
\d\pi(u,v)
=
\begin{pmatrix}
\Sigma & K\\
K^\top & \Lambda
\end{pmatrix}
\]
is positive semidefinite. The Schur complement gives
$R\eqdef\Sigma^{-1/2}K\Lambda^{-1/2}$ with $RR^\top\preceq\Id$, so
$K=\Sigma^{1/2}R\Lambda^{1/2}$ with $\norm{R}_{\mathrm{op}}\leq1$. By duality between the nuclear and operator norms,
\[
\tr(K)
\leq
\sup_{\norm{R}_{\mathrm{op}}\leq1}\tr(\Sigma^{1/2}R\Lambda^{1/2})
=
\norm{\Lambda^{1/2}\Sigma^{1/2}}_*
=
\tr\pa{(\Sigma^{1/2}\Lambda\Sigma^{1/2})^{1/2}}.
\]
For singular matrices, fix $\eta>0$. Adding $\eta\Id$ to the two diagonal blocks preserves block positivity, so the same Schur-complement argument gives
\(\tr(K) \leq \tr\pa{((\Sigma+\eta\Id)^{1/2}(\Lambda+\eta\Id)(\Sigma+\eta\Id)^{1/2})^{1/2}}.\)
Letting $\eta\downarrow0$ and using continuity of the matrix square root gives the same trace bound with $\Sigma,\Lambda$. Hence every admissible coupling has cost at least $\Bb(\Sigma,\Lambda)^2$.

Conversely, choose the centered Gaussian laws $\Gaussian(0,\Sigma)$ and $\Gaussian(0,\Lambda)$. They satisfy the two second-moment constraints, and Proposition~\ref{prop-gaussian-w2-bures} gives
$\Wass_2(\Gaussian(0,\Sigma),\Gaussian(0,\Lambda))^2
=\Bb(\Sigma,\Lambda)^2$, again by continuity in the singular case. This admissible pair attains the lower bound, proving the claim.
\end{proof}
\begin{prop}[Bures metric as a Procrustes metric]\label{prop-bures-procrustes}
Let $\Sigma,\Lambda\in\mathbb S_+^d$, let $d'\geq d$, and fix any
$A_0,B_0\in\RR^{d\times d'}$ such that
$A_0A_0^\top=\Sigma$ and $B_0B_0^\top=\Lambda$. Then
\begin{equation}\label{eq-bures-procrustes}
\Bb(\Sigma,\Lambda)^2
=
\min_{\substack{A,B\in\RR^{d\times d'}\\
AA^\top=\Sigma,\;BB^\top=\Lambda}}
\norm{A-B}_{\mathrm F}^2
=
\min_{R\in\mathrm O(d')}\norm{A_0-B_0R}_{\mathrm F}^2.
\end{equation}
For $d'=d$, one may in particular take $A_0=\Sigma^{1/2}$ and
$B_0=\Lambda^{1/2}$; the second minimum is then over $\mathrm O(d)$.
\end{prop}
\begin{proof}
Two matrices in $\RR^{d\times d'}$ have the same row-Gram matrix if and
only if they differ by right multiplication by an orthogonal matrix. Indeed,
equality of the Gram matrices makes the linear map sending each row of one
matrix to the corresponding row of the other well defined and isometric on
their row spans; this isometry extends to an element of $\mathrm O(d')$.
Therefore every admissible pair has the form $A=A_0U$ and $B=B_0V$ for
some $U,V\in\mathrm O(d')$. Right invariance of the Frobenius norm gives \(\min_{U,V\in\mathrm O(d')}\norm{A_0U-B_0V}_{\mathrm F}^2 = \min_{R\in\mathrm O(d')}\norm{A_0-B_0R}_{\mathrm F}^2,\)
where one may take $R=VU^\top$. Thus the two factor minima in
\eqref{eq-bures-procrustes} coincide and are independent of the chosen
representatives $A_0,B_0$.

Expanding the last objective and applying the orthogonal Procrustes formula,
which follows directly from a singular-value decomposition, yields
\(\min_{R\in\mathrm O(d')}\norm{A_0-B_0R}_{\mathrm F}^2 = \tr(\Sigma)+\tr(\Lambda)-2\norm{A_0^\top B_0}_*.\)
Moreover,
$(A_0^\top B_0)^\top(A_0^\top B_0)=B_0^\top\Sigma B_0$, whose nonzero
eigenvalues are those of
$\Sigma^{1/2}B_0B_0^\top\Sigma^{1/2}
=\Sigma^{1/2}\Lambda\Sigma^{1/2}$.
$\norm{A_0^\top B_0}_*
=\tr((\Sigma^{1/2}\Lambda\Sigma^{1/2})^{1/2})$.
Definition~\ref{def-bures-metric} now identifies the resulting value with
$\Bb(\Sigma,\Lambda)^2$.
\end{proof}
Thus the map $A\mapsto AA^\top$ identifies $\mathbb S_+^d$ with the orbit space $\RR^{d\times d'}/\mathrm O(d')$, and $\Bb$ is the quotient distance induced by the Frobenius metric. This is the matrix counterpart of the quotient transport geometries developed in Section~\ref{sec-quotient-wasserstein-procrustes}.

\begin{prop}[Metric and convexity properties of the Bures term]\label{prop-bures-metric-convex}
The function $\Bb$ is a distance on positive semidefinite covariance matrices. Moreover, $\Bb^2$ is jointly convex: for $t\in[0,1]$,
\(\Bb^2((1-t)\Sigma_0+t\Sigma_1,(1-t)\Lambda_0+t\Lambda_1) \leq (1-t)\Bb^2(\Sigma_0,\Lambda_0)+t\Bb^2(\Sigma_1,\Lambda_1).\)
\end{prop}
\begin{proof}
The free-factor formula in Proposition~\ref{prop-bures-procrustes} is
symmetric and nonnegative. Since its minima are attained, it vanishes if and
only if two factors coincide, which is equivalent to $\Sigma=\Lambda$.
For the triangle inequality, use the square-factor specialization and let
$Q_1,Q_2$ be optimal orthogonal alignments for $(\Sigma,\Lambda)$ and
$(\Lambda,\Gamma)$, respectively.
Since $Q_2Q_1$ is admissible for $(\Sigma,\Gamma)$,
\[
\Bb(\Sigma,\Gamma)
\leq
\norm{\Sigma^{1/2}-\Gamma^{1/2}Q_2Q_1}_{\mathrm F}
\leq
\norm{\Sigma^{1/2}-\Lambda^{1/2}Q_1}_{\mathrm F}
+
\norm{\Lambda^{1/2}-\Gamma^{1/2}Q_2}_{\mathrm F}
=
\Bb(\Sigma,\Lambda)+\Bb(\Lambda,\Gamma).
\]

For convexity, use the free-factor formula in
Proposition~\ref{prop-bures-procrustes}. For $i\in\{0,1\}$, choose optimal
square factors $(U_i,V_i)$ for $(\Sigma_i,\Lambda_i)$ and define the block
factors
\(U_t=[\sqrt{1-t}\,U_0,\sqrt t\,U_1], \qquad V_t=[\sqrt{1-t}\,V_0,\sqrt t\,V_1].\)
Set $\Sigma_t=(1-t)\Sigma_0+t\Sigma_1$ and
$\Lambda_t=(1-t)\Lambda_0+t\Lambda_1$. Then
$U_tU_t^\top=\Sigma_t$ and $V_tV_t^\top=\Lambda_t$. Applying the
free-factor formula with $2d$ columns gives
\[
\Bb^2(\Sigma_t,\Lambda_t)
\leq\norm{U_t-V_t}_{\mathrm F}^2
=(1-t)\norm{U_0-V_0}_{\mathrm F}^2+t\norm{U_1-V_1}_{\mathrm F}^2
=(1-t)\Bb^2(\Sigma_0,\Lambda_0)+t\Bb^2(\Sigma_1,\Lambda_1),
\]
which proves joint convexity.
\end{proof}
\paragraph{Comparison with the Fisher--Rao metric for zero-mean Gaussians.}
\[
\KL(\Sigma\mid\Sigma')
\eqdef
\KL\big(\Gaussian(0,\Sigma)\mid\Gaussian(0,\Sigma')\big)
=
\frac12\left(
\tr((\Sigma')^{-1}\Sigma)-d-\log\det((\Sigma')^{-1}\Sigma)
\right).
\]
If $D=D^\top$ and $\Sigma+\varepsilon D$ remains positive definite, then

\[
\KL(\Sigma+\varepsilon D\mid\Sigma)
=
\frac{\varepsilon^2}{2}\langle Q_\Sigma(D),D\rangle_F+o(\varepsilon^2),
\qquad
Q_\Sigma(D)=\frac12\Sigma^{-1}D\Sigma^{-1},
\]
where $\langle A,B\rangle_F=\tr(A^\top B)$ is the Frobenius pairing. Thus the local quadratic form is

\(g_\Sigma(D,E)=\frac12\tr(\Sigma^{-1}D\Sigma^{-1}E).\)
\(d_{\mathrm{FR}}(\Sigma,\Sigma')^2 = \frac12 \norm{\log(\Sigma^{-1/2}\Sigma'\Sigma^{-1/2})}_F^2,\)
where $\log$ denotes the principal matrix logarithm, and the corresponding geodesic is

\[
\Sigma_t^{\mathrm{FR}}
=
\Sigma^{1/2}
\left(\Sigma^{-1/2}\Sigma'\Sigma^{-1/2}\right)^t
\Sigma^{1/2}.
\]
The contrast with Bures is especially visible near the boundary of the covariance cone. Fisher--Rao is different.

\(\Delta_z=t^2-u^2-v^2,\qquad \langle z,z'\rangle_{\mathrm E}=tt'+uu'+vv', \qquad \langle z,z'\rangle_{\mathrm L}=tt'-uu'-vv'.\)
\begin{align}
\Bb(\Sigma,\Sigma')^2
&=
\sqrt2(t+t')
-2\sqrt{\langle z,z'\rangle_{\mathrm E}
+\sqrt{\Delta_z\Delta_{z'}}},
\label{eq-2x2-bures-cone-distance}\\
d_{\mathrm{FR}}(\Sigma,\Sigma')^2
&=
\frac12\left((\log\lambda_+)^2+(\log\lambda_-)^2\right),
\qquad
\lambda_\pm
=
\frac{\langle z,z'\rangle_{\mathrm L}
\pm\sqrt{\langle z,z'\rangle_{\mathrm L}^2-\Delta_z\Delta_{z'}}}
{\Delta_z},
\label{eq-2x2-fr-cone-distance}
\end{align}
where the Fisher--Rao formula requires $\Delta_z,\Delta_{z'}>0$ and $\lambda_\pm$ are the eigenvalues of $\Sigma^{-1}\Sigma'$.